\scanpage{167}
\mspage{41}

\deleted{En effet} et $C$ et $TC$ sont \inserted{Q-}stables à gauche.

En effet, par le th. 1 du \S\,1 (p. 10) et
son corollaire, $\widetilde C_0$ est la plus petite partie
de $\mathrm{Fl}\unsure{\mathcal M|M}$ contenant $C_0$ et Q-stable à gauche.
Comme $\mathrm{Cof}_W$ \deleted{$\unsure{\subset}$} est Q-stable à g. et contient
$C_0$, on a donc $C=\widetilde C_0\subset\mathrm{Cof}_W$. Donc
$TC\overset{\mathrm{déf}}{=}C\cap W\subset\mathrm{Cof}_W\cap W\overset{\mathrm{déf}}{=}W^{\mathrm{univ.}}$,
\deleted{D'autre part $TC$} donc $TC=C\cap W^{\mathrm{univ}}$. Comme
$W^{\mathrm{univ}}$ est \inserted{lui aussi} Q-stable à g. (cor. précédent), il
en est de $=$ de $C\cap W$, qed.

\uline{Corollaire} 4.4.5. Supposons que \deleted{tout} objet de
$\mathcal M$ soit accessible. Sous les conditions du
corollaire précédent, \deleted{$\unsure{TF=C_*|F=C_*|C_*}$} posant
\[
 TF=C_{0*},=C_*,\qquad \deleted{\unsure{TC=C_*}}\ F=(TC)_*,
\]
le couple $(C,TF)$ est un couple de
Quillen. Pour qu'il en soit de $=$
du couple $(TC,F)$, il suffit qu'il
existe une \uline{petite} partie $TC_0$ de $C\cap W$,
telle que $\widetilde{TC_0}=TC\;(\overset{\mathrm{déf}}{=}C\cap W)$.

La première assertion est, \inserted{on l'a vu tantôt,} le th. 1 du
\S\,1. \deleted{déjà} La deuxième résulte \unsure{de|du} ce $=$
th. On peut préciser dans ce sens :

\uline{Prop. 4.4.\unsure{6|5}} Soit $\mathcal M$ une catégorie stable par petites
$\underrightarrow{\mathrm{lim}}$ filtrantes, par sommes amalgamées, et dont
tous les objets sont accessibles, $W$ un localiseur

% archive sheet number: 166
